% !TEX root = ../main.tex
\begin{titlepage}
\centering
\vspace*{1.2cm}
{\LARGE\bfseries Stacks 패키지 레지스트리를 위한\\[0.3em]
Clarity 스마트 컨트랙트 순위 산정:\\[0.3em]
ERC-20 allowance를 대신하는 포스트컨디션\par}
\vspace{1.4cm}
% 공저자(지도교수 포함 여부)와 교신 저자 이메일은 정하지 않았다.
{\large 권태홍 (Taehong Kwon)\par}
\vspace{0.3cm}
{남아프리카대학교(University of South Africa)\par}
\vspace{1.0cm}
{원고 초안, 2026년 10월 7일\par}
\vfill
\end{titlepage}

\section*{초록}

Stacks에 배포된 Clarity 스마트 컨트랙트는 모두 사람이 읽을 수 있는 소스 형태로 체인에 저장된다. 그런데도 개발자가 이 컨트랙트를 찾아 다시 쓸 수 있는, npm이나 pkg.go.dev 같은 레지스트리는 없다. 이런 레지스트리에는 복제본과 스팸을 사람들이 실제로 의존하는 컨트랙트보다 아래에 두고, 공격자가 가짜 신원으로 싸게 살 수 없는 순위가 필요하다. 선행 연구는 이더리움 지갑의 순위를 ERC-20 전송 위의 PageRank(AWP)나 ERC-20 allowance 위의 PageRank(EndorseRank)로 매겼다. Stacks에는 allowance가 없다. SIP-010 토큰 표준에는 \code{approve}도 \code{transferFrom}도 없다. 사용자는 컨트랙트를 직접 호출해 자기 토큰을 옮기게 하고, 서명한 트랜잭션에 붙인 포스트컨디션으로 자기 계정에서 나갈 수 있는 양을 제한한다. Clarity~4부터는 컨트랙트도 \code{as-contract?}와 \code{restrict-assets?}의 자산 한도로 자기 자산의 유출을 제한한다.

이 논문은 allowance 관계를 이 메커니즘들에 대응시키고, 노드가 컨트랙트뿐인 순위 \PCER를 정의한다. 포스트컨디션 승인을 EndorseRank처럼 지갑·컨트랙트·자산마다 최신 상태 하나로 줄이고, 이것으로 각 지갑이 보증을 여러 컨트랙트에 어떻게 나누는지를 정한다. 지갑이 낸 수수료는 그 보증을 얼마나 무겁게 셀지를 정한다. 컨트랙트 코드의 위임 참조와 의존 참조는 점수를 컨트랙트에서 컨트랙트로 옮긴다. 균등 재시작을 쓰면 컨트랙트 점수 가운데 지갑에서 오는 부분이 지갑 수를 나누어 센 값이 되어 시빌 지갑을 막지 못한다는 것을 보이고, 공격자 컨트랙트 묶음이 모을 수 있는 점수의 상한을 그 묶음의 재시작 몫과 정직한 컨트랙트에서 흘러드는 점수로 나타낸다. 합성 예제로 두 결과를 보인다. 이어서 레지스트리와 명령줄 클라이언트 \cpm을 설명한다. \cpm은 배포된 컨트랙트를 참조하는 일과 그 소스를 복제하는 일을 구분한다. 끝으로 라벨 기간이 열리기 전에 고정한 Stacks 메인넷 평가 계획을 제시한다. 이 평가는 공격자가 컨트랙트를 끌어올리는 데 드는 비용과, 순위가 이후의 채택과 얼마나 맞는지를 개수 기반, 의존성 PageRank, AWP, EndorseRank 기준선과 비교한다.

\textbf{주제어:} Stacks; Clarity; 스마트 컨트랙트; 패키지 레지스트리; PageRank; 시빌 공격 저항; 포스트컨디션; 소프트웨어 공급망

\clearpage
